\chapter{Principio de Acción para Sistemas Covariantes}

\section{Principio de acción de la partícula libre}

El principio de acción afirma que un sistema mecánico posee una cierta
integral, llamada \textbf{principio de acción}, que toma su valor
extremo para el movimiento tal que su variación $\delta S$ es igual a
cero. Determinemos el principio de acción para una partícula libre
relativista, es decir, para una partícula no sometida a una fuerza
externa. Puesto que la física debe ser covariante bajo transformaciones
de Lorentz, el principio de acción debe ser covariante y, como es un
escalar, debe ser \textbf{invariante} bajo transformaciones de Lorentz.

En mecánica clásica el principio de acción está dado por
\[
S = \int_{t_i}^{t_f}L\big(q^i,\dot q^i,t\big)\,dt, \tag{1}
\]
donde $q^i(t)$ depende del tiempo; la integración del Lagrangiano se
realiza en función del tiempo $t$ de un S.R.I. $K$. Sin embargo, el
tiempo también es una coordenada, es decir, es relativo, y por lo tanto
no es posible usar el tiempo de un S.R.I. $K$.

Una partícula que se mueve en el espacio de cuatro dimensiones genera
una línea de mundo.

\begin{center}
\begin{tikzpicture}[scale=1]
  \draw[-Stealth] (0,0) -- (4.4,0) node[right] {$x^1$};
  \draw[-Stealth] (0,0) -- (0,4.4) node[above] {$x^0$};
  \draw[thick,colordef] (1.6,0.4) .. controls (2.3,1.4) and (1.4,2.6) .. (2.2,4.0);
  \node[colordef,right] at (2.3,3.6) {\footnotesize línea de mundo};
  \fill (1.6,0.4) circle (1.5pt) node[left] {$\tau_i$};
  \fill (2.2,4.0) circle (1.5pt) node[left] {$\tau_f$};
\end{tikzpicture}
\end{center}

Por lo tanto, lo que se debe extremizar es la \emph{longitud de la línea
de mundo}: se extremiza su longitud para obtener la ecuación de
movimiento de una partícula libre. Además, la longitud de la línea de
mundo es un escalar invariante de Lorentz. El elemento de línea es
\[
ds^2 = \eta_{\mu\nu}dx^\mu dx^\nu, \tag{2}
\]
y como estudiamos partículas físicas, $ds^2<0$; al extraer la raíz,
$ds=\sqrt{-\eta_{\mu\nu}dx^\mu dx^\nu}$. Como queremos extremizar la
longitud de la línea de mundo, parametrizamos con el tiempo propio,
$x^\mu=x^\mu(\tau)$:
\[
ds = \sqrt{-\eta_{\mu\nu}\frac{dx^\mu}{d\tau}d\tau\frac{dx^\nu}{d\tau}
d\tau} = \sqrt{-\eta_{\mu\nu}u^\mu u^\nu}\,d\tau,
\]
y formamos el funcional
\[
S = \int_{\tau_i}^{\tau_f}\sqrt{-\eta_{\mu\nu}u^\mu u^\nu}\,d\tau .
\tag{3}
\]
Como queremos que el principio de acción tenga unidades de acción,
$[\text{J}\,\text{s}]=[\text{kg}\,\text{m}^2/\text{s}]$, debemos agregar
una constante de dimensionalidad,
\[
S = \lambda\int_{\tau_i}^{\tau_f}\sqrt{-\eta_{\mu\nu}u^\mu u^\nu}\,d\tau .
\tag{4}
\]
Por análisis dimensional, $[\lambda\,ds]=[S]$ implica
$[\lambda]\,\text{m}=\text{kg}\,\text{m}^2/\text{s}$, es decir
\[
[\lambda] = \frac{\text{kg}\,\text{m}}{\text{s}} = [\text{momentum
lineal}], \tag{5}
\]
de modo que la constante debe ser
\[
\lambda = m_0c, \tag{6}
\]
donde $m_0$ es la masa en reposo de la partícula y $c$ la rapidez de la
luz. Entonces
\[
S = m_0c\int_{\tau_i}^{\tau_f}\sqrt{-\eta_{\mu\nu}u^\mu u^\nu}\,d\tau,
\qquad L = m_0c\sqrt{-\eta_{\mu\nu}u^\mu u^\nu}. \tag{7, 8}
\]
Todavía falta un punto importante: el Lagrangiano debe desembocar en la
expresión no relativista de la partícula libre. Escribamos la acción en
el S.R.I. $K$, es decir con respecto al tiempo $t$:
\[
dS = m_0c\sqrt{-\eta_{\mu\nu}dx^\mu dx^\nu} = m_0c\sqrt{c^2dt^2-dx^idx^i}
= m_0c^2\sqrt{1-\frac{v^iv_i}{c^2}}\,dt,
\]
\[
S = m_0c^2\int_{t_i}^{t_f}\sqrt{1-\frac{v^2}{c^2}}\,dt, \qquad
L = m_0c^2\sqrt{1-\frac{v^2}{c^2}} . \tag{9, 10}
\]
Para velocidades bajas o cotidianas, $v\ll c$,
\[
L = m_0c^2\left(1-\frac12\frac{v^2}{c^2}-\frac18\frac{v^4}{c^4}-\dots
\right) = m_0c^2-\frac12m_0v^2-\dots \;\Longrightarrow\; L-m_0c^2
\approx-\frac12m_0v^2 .
\]
El término cinético aparece con un signo menos. Pero el término cinético
siempre debe ser positivo; de lo contrario, cuando se cuantiza el sistema
se presentan irregularidades físicas que no estudiaremos aquí. Entonces
debe cumplirse, para cualquier sistema físico,
\[
T>0. \tag{11}
\]
Por eso agregamos un signo menos a la acción, y finalmente se obtiene el
\textbf{principio de acción de la partícula libre}:
\[
S = -m_0c\int_{\tau_i}^{\tau_f}\sqrt{-\eta_{\mu\nu}u^\mu u^\nu}\,d\tau .
\tag{12}
\]
La causa de que aparezcan tantos signos menos es el uso de la métrica
relativista $\eta_{\mu\nu}=\text{diag}(-1,\delta_{ij})$; si usáramos la
métrica de la física de partículas, el principio tendría signos
positivos.

Recordemos que hay una ligadura que debe satisfacer la cuadrivelocidad,
\[
\eta_{\mu\nu}u^\mu u^\nu = -c^2 . \tag{13}
\]
Si la reemplazamos directamente en (12) se obtiene
$S=-m_0c^2\int d\tau=-m_0c^2\Delta\tau$: la acción se vuelve una
constante y no un principio de acción, y se destruye el principio.
Entonces, \emph{antes} de usar la ligadura (13) se varía primero la
acción y después se reemplaza la ligadura.

\begin{proposicionc}[Ecuación de movimiento de la partícula libre]
Extremizar la acción (12), con $\delta x^\mu(\tau_i)=\delta x^\mu
(\tau_f)=0$, conduce a
\[
\frac{dp^\mu}{d\tau}=0 \qquad\Longrightarrow\qquad p^\mu=\text{constante}.
\tag{14}
\]
\end{proposicionc}

\begin{proof}
Variamos (12):
\[
\delta S = -m_0c\int_{\tau_i}^{\tau_f}\delta\sqrt{-\eta_{\mu\nu}u^\mu
u^\nu}\,d\tau = m_0c\int_{\tau_i}^{\tau_f}\frac{\eta_{\mu\nu}\,
\delta u^\mu\,u^\nu}{\sqrt{-\eta_{\alpha\beta}u^\alpha u^\beta}}\,d\tau
= m_0c\int_{\tau_i}^{\tau_f}\eta_{\mu\nu}\frac{d\,\delta x^\mu}{d\tau}
\frac{u^\nu}{\sqrt{-\eta_{\alpha\beta}u^\alpha u^\beta}}\,d\tau .
\]
Integrando por partes,
\[
\delta S = -m_0c\int_{\tau_i}^{\tau_f}\frac{d}{d\tau}\left(\frac{\eta_{\mu
\nu}u^\nu}{\sqrt{-\eta_{\alpha\beta}u^\alpha u^\beta}}\right)\delta
x^\mu\,d\tau + m_0c\left[\frac{\eta_{\mu\nu}u^\nu}{\sqrt{-\eta_{\alpha
\beta}u^\alpha u^\beta}}\,\delta x^\mu\right]_{\tau_i}^{\tau_f}.
\]
Con $\delta x^\mu(\tau_f)=\delta x^\mu(\tau_i)=0$ se anula el término de
borde, y la condición $\delta S=0$ da
\[
m_0c\,\frac{d}{d\tau}\left(\frac{\eta_{\mu\nu}u^\nu}{\sqrt{-\eta_{\alpha
\beta}u^\alpha u^\beta}}\right)=0 .
\]
Ahora sí introducimos la ligadura (13), $\sqrt{-\eta_{\alpha\beta}
u^\alpha u^\beta}=c$, y se obtiene $m_0\dfrac{d}{d\tau}(\eta_{\mu\nu}
u^\nu)=0$. Como la métrica es constante e invertible, $m_0\dfrac{du^\mu}
{d\tau}=0$, o bien $\dfrac{dp^\mu}{d\tau}=0$. \qedhere
\end{proof}

También se puede extremar la acción escrita en el S.R.I. $K$; el cálculo
es más largo, pero se obtiene lo mismo. Con $S=-m_0c^2\int_{t_i}^{t_f}
\sqrt{1-v^2/c^2}\,dt$,
\[
\delta S = -m_0c^2\int_{t_i}^{t_f}\frac{-v_i\,\delta v^i/c^2}{\sqrt{1-
v^2/c^2}}\,dt = m_0\int_{t_i}^{t_f}\frac{v_i}{\sqrt{1-v^2/c^2}}\,
\frac{d\,\delta x^i}{dt}\,dt .
\]
Integrando por partes y exigiendo $\delta x^i(t_f)=\delta x^i(t_i)=0$,
\[
\delta S = -m_0\int_{t_i}^{t_f}\frac{d}{dt}\left(\frac{v_i}{\sqrt{1-v^2/
c^2}}\right)\delta x^i\,dt = 0 \;\Longrightarrow\; \frac{v^i}{\sqrt{1-v^2
/c^2}} = \gamma v^i=\text{constante}. \tag{15}
\]
La dirección de la velocidad es constante. Tomando el módulo de (15),
$\gamma v=\text{constante}'$, de donde $v=\text{constante}$. Como la
rapidez también es constante, la velocidad es constante, lo que implica
$\gamma=\text{constante}$. Entonces
\[
p^\mu = m_0u^\mu = \big(m_0c\gamma,\ m_0\gamma v^i\big) =
(\text{constante},\text{constante}),
\]
y se recupera (14).

\begin{proposicionc}
La acción de la partícula libre (12) es invariante de Lorentz:
$S'=S$ (16).
\end{proposicionc}

\begin{proof}
Sea la transformación de Lorentz $x'^\mu=\Lambda^\mu_{\ \nu}x^\nu$, tal
que $\eta'_{\mu\nu}=\eta_{\mu\nu}$, $u'^\mu=\Lambda^\mu_{\ \nu}u^\nu$ y
$\tau'=\tau$. Entonces
\[
S' = -m_0c\int_{\tau_i'}^{\tau_f'}\sqrt{-\eta'_{\mu\nu}u'^\mu u'^\nu}\,
d\tau' = -m_0c\int_{\tau_i}^{\tau_f}\sqrt{-\Lambda^\mu_{\ \alpha}\eta_{\mu
\nu}\Lambda^\nu_{\ \beta}u^\alpha u^\beta}\,d\tau .
\]
Pero $\Lambda^\mu_{\ \alpha}\eta_{\mu\nu}\Lambda^\nu_{\ \beta}=
\eta_{\alpha\beta}$, luego $S'=-m_0c\int\sqrt{-\eta_{\alpha\beta}
u^\alpha u^\beta}\,d\tau=S$. \qedhere
\end{proof}

Si la partícula está en presencia de una fuerza conservativa, el
Lagrangiano se escribe
\[
L = -m_0c^2\sqrt{1-\frac{v^2}{c^2}}-U . \tag{17}
\]

\section{Energía y momentum de la partícula libre}

Sea el Lagrangiano de la partícula libre (10) (con el signo definitivo,
$L=-m_0c^2\sqrt{1-v^2/c^2}$), y calculemos el momentum lineal:
\[
p_i = \frac{\partial L}{\partial v^i} = \frac{\partial}{\partial v^i}
\left(-m_0c^2\sqrt{1-\frac{v^2}{c^2}}\right) = \frac{m_0v_i}{\sqrt{1-v^2/
c^2}} = mv_i . \tag{18}
\]
La energía de una partícula libre viene dada por su Hamiltoniano,
\[
E = H = p_iv^i-L = \frac{m_0v^2}{\sqrt{1-v^2/c^2}}+m_0c^2\sqrt{1-
\frac{v^2}{c^2}} = \frac{m_0c^2}{\sqrt{1-v^2/c^2}}\left(\frac{v^2}{c^2}+
1-\frac{v^2}{c^2}\right),
\]
\[
E(v) = \frac{m_0c^2}{\sqrt{1-v^2/c^2}} = mc^2 . \tag{19}
\]
Esta ecuación es de gran importancia porque pone de manifiesto que la
energía de una partícula libre no tiende a cero para velocidad nula,
$v=0$, sino que la energía de una partícula libre en reposo no es nula:
\[
E(0) = m_0c^2 . \tag{20}
\]
Las ecuaciones de movimiento pueden obtenerse del Lagrangiano (17); son
las ecuaciones de Euler--Lagrange,
\[
\frac{d}{dt}\left(\frac{\partial L}{\partial v^i}\right)-\frac{\partial
L}{\partial x^i}=0 \;\Longrightarrow\; \frac{d}{dt}\left(\frac{m_0v_i}
{\sqrt{1-v^2/c^2}}\right)+\frac{\partial U}{\partial x^i}=0
\;\Longrightarrow\; \frac{dp_i}{dt}=F_i . \tag{21}
\]
Debe notarse que el Lagrangiano ya no está dado por $L=T-U$; sin
embargo, su derivada con respecto a la velocidad sigue siendo el
momentum lineal. El momentum canónico en coordenadas generalizadas se
define por $p_i=\partial L/\partial\dot q^i$, de modo que la conexión
entre coordenadas cíclicas y conservación del momentum correspondiente se
mantiene, como en la teoría no relativista. Si $L$ no contiene
explícitamente al tiempo, existe una constante del movimiento,
$H=\dot q^ip_i-L$. Notemos que $L$ dado por (17) no es una función
homogénea de las componentes de la velocidad. Entonces el Hamiltoniano
toma la forma
\[
H = E = \frac{m_0c^2}{\sqrt{1-v^2/c^2}}+U = mc^2+U = (mc^2-m_0c^2)+U+
m_0c^2 = \Delta m\,c^2+U+m_0c^2,
\]
\[
H = E = T+U+m_0c^2 . \tag{22}
\]
La cantidad $H$ se interpreta como la energía total del sistema, que es
por lo tanto una constante del movimiento.

\section{Acción de Maxwell}

\subsection{Las ecuaciones de Maxwell y la conservación de la carga}

Las ecuaciones de Maxwell escritas en forma diferencial son
\[
\begin{aligned}
\vec\nabla\cdot\vec E &= \frac{\rho}{\epsilon_0} &&\text{(ley de Gauss
para el campo eléctrico)}\\
\vec\nabla\times\vec B-\mu_0\epsilon_0\frac{\partial\vec E}{\partial t}
&= \mu_0\vec J &&\text{(ley de Ampère modificada)}\\
\vec\nabla\times\vec E+\frac{\partial\vec B}{\partial t} &= \vec 0
&&\text{(ley de Faraday)}\\
\vec\nabla\cdot\vec B &= 0 &&\text{(ley de Gauss para la inducción
magnética)}
\end{aligned} \tag{23}
\]
donde $\vec E$ es el campo eléctrico, $\vec B$ la inducción magnética,
$\rho$ la densidad de carga eléctrica, $\vec J$ la densidad de corriente
por área transversal, $\epsilon_0=8.854\times10^{-12}\ \text{s}^2\text{C}^2/
(\text{kg}\,\text{m}^3)$ la permitividad eléctrica del vacío y
$\mu_0=4\pi\times10^{-7}\ \text{kg}\,\text{m}/\text{C}^2$ la
permeabilidad magnética del vacío. Nos quedaremos con la forma
diferencial (las integrales solo se pueden resolver si el fenómeno posee
una alta simetría).

Las ecuaciones de Maxwell involucran seis incógnitas, las tres
componentes de $\vec E$ y las tres de $\vec B$, que se pueden resolver si
son conocidas $\rho$ y $\vec J$. Sin embargo, las ecuaciones son ocho
(dos vectoriales y dos escalares). Se podría pensar que el sistema está
sobredimensionado, pero no es así: las ecuaciones que sobran, en este
caso dos, relacionan el campo eléctrico y el magnético cuando hay
evolución espacio-temporal.

Queremos poner las ecuaciones de Maxwell en forma covariante de Lorentz,
es decir, que sean tensores de Lorentz. Comencemos con un análisis
dimensional (¡esto ayuda mucho!):
\[
[E]=\left[\frac Fq\right]=\frac{\text{kg}\,\text{m}}{\text{s}^2\,
\text{C}}, \quad [B]=\left[\frac{F}{qv}\right]=\frac{\text{kg}}{\text{s}\,
\text{C}}, \quad [\rho]=\frac{\text{C}}{\text{m}^3}, \quad [J]=\frac{
\text{A}}{\text{m}^2}.
\]
El campo eléctrico y la inducción magnética no poseen las mismas
unidades: $[E/B]=\text{m/s}$, hay una rapidez que las relaciona y, como
veremos, es la rapidez de la luz.

\begin{proposicionc}[Ecuación de continuidad]
Las ecuaciones de Maxwell implican la conservación de la carga:
\[
\frac{\partial\rho}{\partial t}+\vec\nabla\cdot\vec J=0 . \tag{26}
\]
\end{proposicionc}

\begin{proof}
Derivando con respecto al tiempo la ley de Gauss para el campo
eléctrico,
\[
\epsilon_0\vec\nabla\cdot\frac{\partial\vec E}{\partial t} = \frac{
\partial\rho}{\partial t}. \tag{24}
\]
Tomando la divergencia de la ley de Ampère modificada y usando
$\vec\nabla\cdot(\vec\nabla\times\vec B)=0$,
\[
-\epsilon_0\vec\nabla\cdot\frac{\partial\vec E}{\partial t} =
\vec\nabla\cdot\vec J . \tag{25}
\]
Sumando (24) y (25) se obtiene (26). \qedhere
\end{proof}

Para entender la ecuación de continuidad integrémosla en una región de
volumen $V$ que encierra a la carga eléctrica $Q$:
\[
\frac{d}{dt}\int_V\rho\,dV+\oint_S\vec J\cdot d\vec S=0
\;\Longrightarrow\; \frac{dQ}{dt} = -\oint_S\vec J\cdot d\vec S, \qquad
S=\partial V. \tag{27}
\]
La corriente que cruza la superficie cerrada $S$ representa la rapidez
con que la carga abandona el volumen encerrado. Una de las leyes básicas
de la física es que la carga es indestructible: nunca se destruye ni se
crea, es decir, \textbf{la carga se conserva}. Si hay una corriente neta
saliendo de una superficie cerrada, la cantidad de carga en el interior
debe disminuir en la misma cantidad que sale; si no hay aporte ni retiro
neto de carga, $dQ/dt=0$ y $\oint_S\vec J\cdot d\vec S=0$.

Escribamos la ecuación de continuidad en forma covariante:
\[
\frac1c\frac{\partial}{\partial t}(c\rho)+\vec\nabla\cdot\vec J=0
\;\Longrightarrow\; \frac{\partial(c\rho)}{\partial(ct)}+\vec\nabla\cdot
\vec J=0. \tag{28}
\]
La cantidad $c\rho$ tiene unidades $[c\rho]=\text{A/m}^2$, las mismas que
$\vec J$; definimos $J^0=c\rho$ (29), y (28) se reescribe
\[
\frac{\partial J^0}{\partial x^0}+\frac{\partial J^i}{\partial x^i}=0
\;\Longrightarrow\; \frac{\partial J^\mu}{\partial x^\mu}=0, \tag{30}
\]
donde se introdujo la \textbf{cuadridensidad de corriente}, con el índice
arriba,
\[
J^\mu = \big(J^0,J^i\big) = \big(c\rho,J^i\big). \tag{31}
\]

\subsection{Las ecuaciones con fuentes en forma covariante}

Introduciendo (29) en la ley de Gauss, $\vec\nabla\cdot\vec E=J^0/(c
\epsilon_0)$ (32), y usando $\mu_0=1/(c^2\epsilon_0)$ (33), al multiplicar
la ley de Ampère modificada por $c$ se tiene
\[
\vec\nabla\times(c\vec B)-\frac1c\frac{\partial\vec E}{\partial t} =
\frac{1}{c\epsilon_0}\vec J . \tag{35}
\]
Notemos que aparece $c\vec B$, cuya dimensionalidad es
$[cB]=\text{kg}\,\text{m}/(\text{s}^2\text{C})$, la misma del campo
eléctrico. Definimos el \textbf{campo magnético}
\[
\vec H = c\vec B . \tag{36}
\]
Con (36) las ecuaciones de Maxwell se escriben
\[
\begin{aligned}
\vec\nabla\cdot\vec E &= \frac{J^0}{c\epsilon_0}, &
\vec\nabla\times\vec H-\frac{\partial\vec E}{\partial x^0} &=
\frac{1}{c\epsilon_0}\vec J,\\
\vec\nabla\times\vec E+\frac{\partial\vec H}{\partial x^0} &= \vec0, &
\vec\nabla\cdot\vec H &= 0,
\end{aligned} \tag{37}
\]
y en notación de índices,
\[
\begin{aligned}
\frac{\partial E_i}{\partial x^i} &= \frac{J^0}{c\epsilon_0}, &
\epsilon_{ijk}\frac{\partial H_k}{\partial x^j}-\frac{\partial E_i}
{\partial x^0} &= \frac{1}{c\epsilon_0}J^i,\\
\epsilon_{ijk}\frac{\partial E_k}{\partial x^j}+\frac{\partial H_i}
{\partial x^0} &= 0, & \frac{\partial H_i}{\partial x^i} &= 0.
\end{aligned} \tag{38}
\]
Para ponerlas en forma covariante necesitamos introducir el campo
eléctrico y el magnético en una sola entidad. Como suman seis
componentes, buscamos un tensor en cuatro dimensiones con seis
componentes independientes: un tensor de orden $(2,0)$ antisimétrico,
\[
F^{\mu\nu}=-F^{\nu\mu}, \tag{39}
\]
que posee $\frac12\times4\times(4-1)=6$ componentes independientes.
Definimos
\[
F^{0i}=-E_i, \qquad F^{ij}=-\epsilon_{ijk}H_k, \tag{40}
\]
de modo que
\[
F^{\mu\nu} = \begin{pmatrix} 0 & F^{0j}\\ F^{i0} & F^{ij}\end{pmatrix}
= \begin{pmatrix} 0 & -E_1 & -E_2 & -E_3\\ E_1 & 0 & -H_3 & H_2\\
E_2 & H_3 & 0 & -H_1\\ E_3 & -H_2 & H_1 & 0\end{pmatrix}. \tag{41, 42}
\]
El tensor $F^{\mu\nu}$ se llama \textbf{tensor del campo
electromagnético} o \textbf{tensor de Faraday}.

Introduzcámoslo primero en las ecuaciones con materia. La ley de Gauss
$\partial E_i/\partial x^i=J^0/(c\epsilon_0)$ se escribe, con
$E_i=-F^{0i}$ y $F^{00}=0$,
\[
\frac{\partial F^{0i}}{\partial x^i} = -\frac{1}{c\epsilon_0}J^0
\;\Longrightarrow\; \frac{\partial F^{0\nu}}{\partial x^\nu} =
-\frac{1}{c\epsilon_0}J^0. \tag{43, 44}
\]
En la ley de Ampère modificada, usando (40), $\epsilon_{ijk}H_k=-F^{ij}$
y $E_i=-F^{0i}=F^{i0}$:
\[
-\frac{\partial F^{ij}}{\partial x^j}-\frac{\partial F^{i0}}{\partial x^0}
= \frac{1}{c\epsilon_0}J^i \;\Longrightarrow\; \frac{\partial F^{i\nu}}
{\partial x^\nu} = -\frac{1}{c\epsilon_0}J^i . \tag{45}
\]
Resumiendo (44) y (45):
\[
\frac{\partial F^{0\nu}}{\partial x^\nu} = -\frac{1}{c\epsilon_0}J^0,
\qquad \frac{\partial F^{i\nu}}{\partial x^\nu} = -\frac{1}{c\epsilon_0}
J^i \tag{46}
\]
\[
\Longrightarrow\quad \boxed{\ \frac{\partial F^{\mu\nu}}{\partial x^\nu}
= -\frac{1}{c\epsilon_0}J^\mu, \qquad \mu,\nu=0,1,2,3\ } \tag{47}
\]
que se llama \textbf{primera forma tensorial de Maxwell}.

De la ecuación (47) se puede obtener la ecuación de continuidad (30):
contrayendo con $\partial/\partial x^\mu$ y usando la antisimetría (39),
$\partial_\mu\partial_\nu F^{\mu\nu}=0$, de modo que
$\partial_\mu J^\mu=0$.

\subsection{Las ecuaciones sin fuentes en forma covariante}

La ecuación (47) no incluye la ley de Faraday ni la ley de Gauss para el
campo magnético,
\[
\vec\nabla\times\vec E+\frac{\partial\vec B}{\partial t}=\vec0, \qquad
\vec\nabla\cdot\vec B=0, \tag{48}
\]
que son cuatro ecuaciones (tres componentes de Faraday y una de Gauss
magnética). Debemos encontrar una ecuación tensorial con cuatro
componentes independientes. Escribamos la ley de Faraday con el tensor
de Faraday: $\epsilon_{ijk}\partial_jE_k+\partial_0H_i=0$ con
$E_k=-F^{0k}$. De $F^{ij}=-\epsilon_{ijk}H_k$ se sigue
\[
H_i = -\frac12\epsilon_{ijk}F^{jk}. \tag{49}
\]
Entonces $-\epsilon_{ijk}\partial_jF^{0k}-\frac12\epsilon_{ijk}\partial_0
F^{jk}=0$; contrayendo con $-\epsilon_{nmi}$ y usando
$\epsilon_{nmi}\epsilon_{ijk}=\delta_{nj}\delta_{mk}-\delta_{nk}
\delta_{mj}$,
\[
\frac{\partial F^{0m}}{\partial x^n}+\frac{\partial F^{n0}}{\partial
x^m}+\frac{\partial F^{nm}}{\partial x^0}=0. \tag{50}
\]
La ley de Gauss para la inducción magnética, $\partial_iH_i=0$, con (49)
se escribe $\epsilon_{ijk}\partial_iF^{jk}=0$. Resumiendo,
\[
\frac{\partial F^{0m}}{\partial x^n}+\frac{\partial F^{n0}}{\partial
x^m}+\frac{\partial F^{nm}}{\partial x^0}=0, \qquad
\epsilon_{ijk}\frac{\partial F^{jk}}{\partial x^i}=0. \tag{51}
\]
La segunda expresión se puede escribir de forma análoga a la primera: como
$\epsilon_{ijk}=\epsilon_{kij}=\epsilon_{jki}$ (permutación cíclica),
\[
\frac13\epsilon_{ijk}\left(\frac{\partial F^{jk}}{\partial x^i}+
\frac{\partial F^{ij}}{\partial x^k}+\frac{\partial F^{ki}}{\partial
x^j}\right)=0. \tag{52}
\]

\begin{definicion}[Tensor de Faraday con índices abajo]
Se define el tensor de Faraday $F_{\mu\nu}$, de orden $(0,2)$, bajando
los índices con la métrica de Minkowski,
\[
F_{\mu\nu} = \eta_{\mu\alpha}\eta_{\nu\beta}F^{\alpha\beta}. \tag{54}
\]
\end{definicion}

\begin{proposicionc}
Si el tensor $F^{\mu\nu}$ es antisimétrico, entonces $F_{\mu\nu}$
también lo es: $F_{\mu\nu}=-F_{\nu\mu}$ (55).
\end{proposicionc}

\begin{proof}
$F_{\mu\nu}=\eta_{\mu\alpha}\eta_{\nu\beta}F^{\alpha\beta}=\eta_{\mu\beta}
\eta_{\nu\alpha}F^{\beta\alpha}=-\eta_{\nu\alpha}\eta_{\mu\beta}
F^{\alpha\beta}=-F_{\nu\mu}$, donde se renombraron los índices mudos y se
usó (39). \qedhere
\end{proof}

Descomponiendo en componentes temporales y espaciales:
$F_{0i}=\eta_{00}\eta_{ij}F^{0j}=-F^{0i}=E_i$ y
$F_{ij}=F^{ij}=-\epsilon_{ijk}H_k$. Resumiendo,
\[
F_{0i}=E_i=-F^{0i}, \qquad F_{ij}=-\epsilon_{ijk}H_k=F^{ij}, \tag{56}
\]
\[
F_{\mu\nu}=\begin{pmatrix}0&E_1&E_2&E_3\\-E_1&0&-H_3&H_2\\-E_2&H_3&0&-H_1
\\-E_3&-H_2&H_1&0\end{pmatrix}. \tag{57}
\]
Usando (56), las expresiones (53) se reescriben
\[
\frac{\partial F_{m0}}{\partial x^n}+\frac{\partial F_{0n}}{\partial
x^m}+\frac{\partial F_{nm}}{\partial x^0}=0, \qquad
\epsilon_{ijk}\left(\frac{\partial F_{jk}}{\partial x^i}+\frac{\partial
F_{ij}}{\partial x^k}+\frac{\partial F_{ki}}{\partial x^j}\right)=0.
\tag{58}
\]
Por simplicidad definimos el tensor $\partial_\mu:=\partial/\partial
x^\mu$, de orden $(0,1)$ (59), y (58) se reescribe
\[
\partial_0F_{nm}+\partial_nF_{m0}+\partial_mF_{0n}=0, \qquad
\epsilon_{ijk}\big(\partial_iF_{jk}+\partial_kF_{ij}+\partial_jF_{ki}
\big)=0. \tag{60}
\]
Hay una permutación cíclica en los índices $0nm$ e $ijk$, típica cuando
se simetriza o antisimetriza un tensor de tres índices. Antisimetricemos
completamente $\partial_\mu F_{\nu\lambda}$:
\[
S_{\mu\nu\lambda} := \partial_{[\mu}F_{\nu\lambda]} = \frac13\big(
\partial_\mu F_{\nu\lambda}+\partial_\lambda F_{\mu\nu}+\partial_\nu
F_{\lambda\mu}\big), \tag{61}
\]
donde se usó la antisimetría (55). El tensor $S_{\mu\nu\lambda}$ posee
$\binom43=4$ componentes independientes, que se descomponen en un vector
tridimensional y un escalar. Obtengamos las ecuaciones sin fuentes.

\begin{itemize}
\item Con $\mu=0$, $\nu=i$, $\lambda=j$:
\[
S_{0ij}=\frac13\big(\partial_0F_{ij}+\partial_jF_{0i}-\partial_iF_{0j}
\big)=-\frac13\epsilon_{ijk}\big(\partial_0H_k+\epsilon_{kln}\partial_l
E_n\big), \tag{63}
\]
que es la ley de Faraday; por lo tanto $S_{0ij}=0$ (62).
\item Con $\mu=i$, $\nu=j$, $\lambda=k$: $S_{ijk}$ es un tensor
tridimensional totalmente antisimétrico, con $\binom33=1$ componente
independiente, luego es proporcional al símbolo de Levi-Civita,
$S_{ijk}=\alpha\epsilon_{ijk}$ (64). Contrayendo con $\epsilon_{ijk}$,
$\epsilon_{ijk}S_{ijk}=\epsilon_{ijk}\partial_iF_{jk}=-2\partial_nH_n=
3!\,\alpha$, de donde $\alpha=-\frac13\partial_nH_n$ (65) y
\[
S_{ijk}=-\frac13\epsilon_{ijk}\,\partial_nH_n . \tag{66}
\]
Como $\partial_nH_n=0$, $S_{ijk}=0$ (67).
\end{itemize}

Resumiendo,
\[
\left.\begin{aligned}S_{0ij}&=0\\S_{ijk}&=0\end{aligned}\right\}
\Longrightarrow S_{\mu\nu\lambda}=0\;\Longrightarrow\;
\partial_\mu F_{\nu\lambda}+\partial_\lambda F_{\mu\nu}+\partial_\nu
F_{\lambda\mu}=0, \tag{68}
\]
llamada \textbf{segunda forma tensorial de Maxwell}. Las ecuaciones de
Maxwell en forma covariante son entonces
\[
\partial_\nu F^{\mu\nu}=-\frac{1}{c\epsilon_0}J^\mu, \qquad
\partial_\mu F_{\nu\lambda}+\partial_\lambda F_{\mu\nu}+\partial_\nu
F_{\lambda\mu}=0. \tag{69}
\]
Si se contrae la primera con $\partial_\mu$, la antisimetría de
$F^{\mu\nu}$ da $\partial_\mu J^\mu=0$: se obtiene la ecuación de
continuidad.

\begin{definicion}[Dual de Hodge]
Se define el \textbf{tensor de Maxwell dual de Hodge} $^*F^{\mu\nu}$
asociado a $F^{\mu\nu}$ como
\[
{}^*F^{\mu\nu} = \frac12\epsilon^{\mu\nu\alpha\beta}F_{\alpha\beta}.
\tag{70}
\]
\end{definicion}

El tensor $^*F^{\mu\nu}$ está relacionado con los campos eléctrico y
magnético en la forma
\[
{}^*F^{0i}=H_i, \qquad {}^*F^{ij}=-\epsilon_{ijk}E_k, \tag{71}
\]
\[
{}^*F^{\mu\nu}=\begin{pmatrix}0&H_1&H_2&H_3\\-H_1&0&-E_3&E_2\\-H_2&E_3&0&
-E_1\\-H_3&-E_2&E_1&0\end{pmatrix}. \tag{72}
\]
La operación dual de Hodge intercambia el campo eléctrico por el
magnético y viceversa. Contrayendo la segunda ecuación de (69) con
$\epsilon^{\alpha\mu\nu\lambda}$ y usando que los tres términos son
iguales por permutación cíclica,
\[
3\,\epsilon^{\alpha\mu\nu\lambda}\partial_\mu F_{\nu\lambda}=0
\;\Longrightarrow\; \partial_\mu\,{}^*F^{\alpha\mu}=0, \tag{73, 74}
\]
y las ecuaciones tensoriales de Maxwell se reescriben, de forma más
estética,
\[
\partial_\nu F^{\mu\nu}=-\frac{1}{c\epsilon_0}J^\mu, \qquad
\partial_\nu\,{}^*F^{\mu\nu}=0. \tag{75}
\]

\subsection{Covariancia de Lorentz de las ecuaciones de Maxwell}

Se han escrito las ecuaciones de Maxwell en notación de índices; ahora
debemos ver si son covariantes de Lorentz. Una transformación de
Lorentz es $x'^\mu=\Lambda^\mu_{\ \nu}x^\nu$ (76), tal que
\[
\frac{\partial}{\partial x'^\mu}=(\Lambda^{-1})^\nu_{\ \mu}\frac{
\partial}{\partial x^\nu}\;\Longrightarrow\;\partial'_\mu=(\Lambda^{-1})^
\nu_{\ \mu}\partial_\nu. \tag{77}
\]
La cuadridensidad de corriente $J^\mu$ contiene a la tricorriente $\vec
J$, que es un vector y por lo tanto debe satisfacer el principio de
covariancia: $J^\mu$ es un tensor de Lorentz,
\[
J'^\mu=\Lambda^\mu_{\ \nu}J^\nu . \tag{78}
\]
Como consecuencia, de la primera forma tensorial de Maxwell, la cantidad
$\partial_\nu F^{\mu\nu}$ debe ser un tensor de Lorentz. Por lo tanto:

\begin{definicion}
El tensor de Faraday $F^{\mu\nu}$ es un tensor de Lorentz de orden
$(2,0)$:
\[
F'^{\mu\nu}=\Lambda^\mu_{\ \alpha}\Lambda^\nu_{\ \beta}F^{\alpha\beta}.
\tag{79}
\]
\end{definicion}

La ecuación (79) informa cómo transforma el campo electromagnético bajo
una transformación de Lorentz. Descomponiéndola en campo eléctrico y
magnético se encuentra
\[
\begin{aligned}
E'^i &= \gamma\big((\delta^i_j-\hat n^i\hat n_j)E^j-\epsilon^{ijk}\beta_j
H_k\big)+\hat n^i\hat n_jE^j,\\
H'^i &= \gamma\big((\delta^i_j-\hat n^i\hat n_j)H^j+\epsilon^{ijk}
\beta_jE_k\big)+\hat n^i\hat n_jH^j,
\end{aligned} \tag{80}
\]
o, en forma vectorial, separando componentes perpendiculares y paralelas
a $\hat n$,
\[
\vec E\,'=\gamma\big(\vec E_\perp-\vec\beta\times\vec H\big)+\vec E_\parallel,
\qquad \vec H\,'=\gamma\big(\vec H_\perp+\vec\beta\times\vec E\big)+
\vec H_\parallel. \tag{81}
\]
(Aquí $\vec\beta$ corresponde al boost $\Lambda(\beta)$ del Capítulo~7;
si se invierte el sentido de la velocidad relativa, cambia el signo de
$\vec\beta$.) De forma análoga, $F_{\mu\nu}$ es un tensor de orden
$(0,2)$,
\[
F'_{\mu\nu}=(\Lambda^{-1})^\alpha_{\ \mu}(\Lambda^{-1})^\beta_{\ \nu}
F_{\alpha\beta}, \tag{82}
\]
y al reemplazar las expresiones de $F_{\mu\nu}$ y $(\Lambda^{-1})^\alpha_{
\ \mu}$ se recupera (81). Las ecuaciones (69) transforman como tensores
de Lorentz:
\[
\begin{aligned}
\partial'_\lambda F'^{\mu\lambda}+\frac{1}{c\epsilon_0}J'^\mu &=
\Lambda^\mu_{\ \nu}\left(\partial_\alpha F^{\nu\alpha}+\frac{1}{c
\epsilon_0}J^\nu\right),\\
\partial'_\mu F'_{\nu\lambda}+\partial'_\lambda F'_{\mu\nu}+\partial'_\nu
F'_{\lambda\mu} &= (\Lambda^{-1})^\alpha_{\ \mu}(\Lambda^{-1})^\beta_{\
\nu}(\Lambda^{-1})^\gamma_{\ \lambda}\big(\partial_\alpha F_{\beta\gamma}+
\partial_\gamma F_{\alpha\beta}+\partial_\beta F_{\gamma\alpha}\big),
\end{aligned} \tag{83}
\]
lo que era de esperar, pues se respetan los postulados de Einstein.

\subsection{El cuadripotencial electromagnético}

Lo que sigue es obtener el principio de acción para las ecuaciones de
Maxwell,
\[
S_{\text{Maxwell}}=\int_{M^4}\mathscr{L}(\vec E,\vec B)\,d^4x, \tag{84}
\]
tal que la variación con respecto a los campos proporcione las
ecuaciones de campo de Maxwell. Hay que resolver dos problemas:
\begin{itemize}
\item Hay dos ecuaciones de campo, (69), y se podría pensar que son
independientes porque hay dos campos que variar, $\vec E$ y $\vec B$.
Pero el tensor de Faraday expresa que el campo eléctrico y el magnético
son partes de una única entidad física, $F^{\mu\nu}$, por lo que lo
correcto sería variar con respecto a este tensor y se obtendría una sola
ecuación de campo, lo que contradice (69).
\item Las ecuaciones (69) son ecuaciones diferenciales parciales de
primer orden en $\vec E$ y $\vec B$, lo que complica obtener un
principio de acción: si el principio de acción contiene primeras
derivadas, el proceso variacional produce ecuaciones de campo con
derivadas de segundo orden.
\end{itemize}
Ambos inconvenientes se resuelven convirtiendo las ecuaciones (69) en
ecuaciones de segundo orden, resolviendo la ley de Faraday y la de Gauss
para el campo magnético. Partamos con la ley de Gauss magnética:
\[
\vec\nabla\cdot\vec H=0\;\Longrightarrow\;\vec H=-\vec\nabla\times\vec A,
\tag{85}
\]
y llamaremos al vector $\vec A$ \textbf{potencial vectorial magnético}
(abuso del lenguaje). Sus unidades son $[A]=[H]/[\nabla]=\text{kg}\,
\text{m}^2/(\text{s}^2\text{C})$ (86). Históricamente el nombre se asoció
a la inducción $\vec B=-\vec\nabla\times\vec A_B$, y como $\vec H=c\vec B$
se tiene $\vec A_H=c\vec A_B$, con $[A_B]=\text{kg}\,\text{m}/(\text{s}\,
\text{C})$ (87). En lo que sigue llamaremos potencial vectorial magnético
a $\vec A$ relacionado con $\vec H$.

Introduciendo (85) en la ley de Faraday,
\[
\vec\nabla\times\vec E+\frac{\partial}{\partial x^0}\big(-\vec\nabla\times
\vec A\big)=0\;\Longrightarrow\;\vec\nabla\times\left(\vec E-\frac{
\partial\vec A}{\partial x^0}\right)=0\;\Longrightarrow\;\vec E-\frac{
\partial\vec A}{\partial x^0}=-\vec\nabla\phi, \tag{88}
\]
donde $\phi$ es el \textbf{potencial eléctrico}. Resumiendo,
\[
\vec E=\frac{\partial\vec A}{\partial x^0}-\vec\nabla\phi, \qquad
\vec H=-\vec\nabla\times\vec A, \tag{89}
\]
o, en notación de índices,
\[
E_i=\partial_0A_i-\partial_i\phi, \qquad H_i=-\epsilon_{ijk}\partial_jA_k.
\tag{90}
\]
Reemplazando (90) en (56):
\[
F_{0i}=E_i=\partial_0A_i-\partial_i\phi, \qquad F_{ij}=-\epsilon_{ijk}H_k=
\epsilon_{ijk}\epsilon_{lnk}\partial_lA_n=\partial_iA_j-\partial_jA_i.
\tag{91}
\]
Notemos que en (91) aparece $\partial A_i/\partial t$, la variación del
potencial vectorial con el tiempo, pero no aparece la variación de
$\phi$ con el tiempo: $\phi$ no posee una velocidad y, como veremos, no
hay término cinético para $\phi$. Esta situación es muy importante en
teoría de gauge y en el programa de Dirac.

\begin{definicion}[Cuadripotencial electromagnético]
\[
A_\mu = \big(A_0,A_i\big)=\big(\phi,A_i\big), \qquad \mu=0,1,2,3.
\tag{92}
\]
\end{definicion}

Con esta definición, las expresiones (91) se reducen a
\[
F_{\mu\nu}=\partial_\mu A_\nu-\partial_\nu A_\mu . \tag{93}
\]
¿Es el cuadripotencial un tensor de Lorentz de orden $(0,1)$?
Veamos: bajo una transformación de Lorentz, $F_{\mu\nu}$ transforma según
(82); introduciendo (93),
\[
\partial'_\mu A'_\nu-\partial'_\nu A'_\mu=(\Lambda^{-1})^\alpha_{\ \mu}
(\Lambda^{-1})^\beta_{\ \nu}\big(\partial_\alpha A_\beta-\partial_\beta
A_\alpha\big),
\]
y comparando se encuentra
\[
A'_\mu=(\Lambda^{-1})^\alpha_{\ \mu}A_\alpha, \tag{94}
\]
es decir, $A_\mu$ es un tensor de Lorentz de orden $(0,1)$.

\begin{definicion}[Tensor de Faraday]
El tensor de Faraday, de orden $(0,2)$, se escribe en términos del
cuadripotencial como
\[
F_{\mu\nu}=\partial_\mu A_\nu-\partial_\nu A_\mu . \tag{95}
\]
Subiendo los índices con la métrica inversa,
$F^{\mu\nu}=\partial^\mu A^\nu-\partial^\nu A^\mu$ (96).
\end{definicion}

Al reemplazar (96) en la primera ecuación de (69) se encuentra la
ecuación de campo para el cuadripotencial:
\[
\Box A^\mu-\partial^\mu\big(\partial_\nu A^\nu\big)=\frac{1}{c\epsilon_0}
J^\mu, \tag{97}
\]
donde
\[
\Box=\eta^{\mu\nu}\partial_\mu\partial_\nu=\nabla^2-\frac{1}{c^2}
\frac{\partial^2}{\partial t^2} \tag{98}
\]
es el \textbf{operador d'Alembertiano}. La ecuación (97) no es una
ecuación de onda, pues aparece $\partial^\mu(\partial_\nu A^\nu)$; son
cuatro ecuaciones para las cuatro componentes de $A^\mu$, de modo que,
conocido $J^\mu$, podríamos conocer unívocamente el cuadripotencial. Sin
embargo, como veremos, la solución de (97) no es única, debido a una
simetría escondida en la teoría electromagnética.

Reemplacemos (95) en la segunda ecuación de (69):
\[
\begin{aligned}
\partial_\mu F_{\nu\lambda}+\partial_\lambda F_{\mu\nu}+\partial_\nu
F_{\lambda\mu} &= \partial_\mu(\partial_\nu A_\lambda-\partial_\lambda
A_\nu)+\partial_\lambda(\partial_\mu A_\nu-\partial_\nu A_\mu)+\partial_\nu
(\partial_\lambda A_\mu-\partial_\mu A_\lambda)\\
&= (\partial_\mu\partial_\nu-\partial_\nu\partial_\mu)A_\lambda+
(\partial_\lambda\partial_\mu-\partial_\mu\partial_\lambda)A_\nu+
(\partial_\nu\partial_\lambda-\partial_\lambda\partial_\nu)A_\mu \equiv 0.
\end{aligned} \tag{99}
\]
La segunda forma de (69) se anula idénticamente si el tensor de Faraday
se escribe como (95): proviene de resolver la ley de Gauss magnética y la
ley de Faraday. La expresión (95) elimina la segunda ecuación de (69),
dejando la primera, lo que permite obtener un principio de acción.

La ecuación (97) acopla el campo eléctrico y el magnético (contiene la
ley de Ampère modificada), pero se pueden desacoplar. Bajando el índice
en (97), $\Box A_\nu-\partial_\nu(\partial_\lambda A^\lambda)=\dfrac{1}
{c\epsilon_0}J_\nu$ (100). Derivando con respecto a $x^\mu$,
\[
\Box(\partial_\mu A_\nu)-\partial_\mu\partial_\nu(\partial_\lambda
A^\lambda)=\frac{1}{c\epsilon_0}\partial_\mu J_\nu, \tag{101}
\]
e intercambiando $\mu\leftrightarrow\nu$ (102) y restando,
\[
\Box(\partial_\mu A_\nu-\partial_\nu A_\mu)=\frac{1}{c\epsilon_0}
(\partial_\mu J_\nu-\partial_\nu J_\mu)\;\Longrightarrow\;\Box F_{\mu\nu}
=\frac{1}{c\epsilon_0}\big(\partial_\mu J_\nu-\partial_\nu J_\mu\big),
\tag{103}
\]
donde se usó que las derivadas parciales conmutan. Haciendo $\mu=0$,
$\nu=i$ en (103), con $F_{0i}=E_i$, $J_0=-c\rho$ y $J_i=J^i$,
\[
\Box\vec E=\frac{1}{\epsilon_0}\left(\vec\nabla\rho+\frac{1}{c^2}
\frac{\partial\vec J}{\partial t}\right), \tag{104}
\]
y con $\mu=i$, $\nu=j$ (contrayendo con $\epsilon_{ijn}$),
\[
\Box\vec H=-\frac{1}{c\epsilon_0}\vec\nabla\times\vec J . \tag{105}
\]
A partir de (97) se han desacoplado las ecuaciones para los campos
eléctrico y magnético, obteniendo ecuaciones de onda con fuentes del
tipo
\[
\Box\Psi(x^\mu)=-Q(x^\mu), \tag{106}
\]
donde $Q(x^\mu)$ es la fuente. La expresión (103) es un tensor de Lorentz
de orden $(0,2)$:
\[
\Box'F'_{\mu\nu}-\frac{1}{c\epsilon_0}\big(\partial'_\mu J'_\nu-
\partial'_\nu J'_\mu\big)=(\Lambda^{-1})^\alpha_{\ \mu}(\Lambda^{-1})^\beta_
{\ \nu}\left(\Box F_{\alpha\beta}-\frac{1}{c\epsilon_0}(\partial_\alpha
J_\beta-\partial_\beta J_\alpha)\right). \tag{107}
\]

\subsection{El principio de acción de Maxwell}

Para obtener el principio de acción del campo electromagnético recurrimos
a ingeniería inversa. Contraemos la primera ecuación de (69) con
$\delta A_\mu$:
\[
\partial_\nu F^{\mu\nu}\delta A_\mu+\frac{1}{c\epsilon_0}J^\mu\delta
A_\mu=0. \tag{109}
\]
La cuadridensidad de corriente $J^\mu(x^\nu)$ es la fuente del
cuadripotencial, por lo tanto no depende de $A_\mu$, sino solo de las
coordenadas: $\delta J^\mu=0$ (110). Derivando por partes,
\[
-F^{\mu\nu}\partial_\nu\delta A_\mu+\delta\left(\frac{1}{c\epsilon_0}
J^\mu A_\mu\right)+\partial_\nu\big(F^{\mu\nu}\delta A_\mu\big)=0.
\]
Como $F^{\mu\nu}$ es antisimétrico, $F^{\mu\nu}\partial_\nu\delta A_\mu=
\frac12F^{\mu\nu}\delta(\partial_\nu A_\mu-\partial_\mu A_\nu)=
\frac12F^{\mu\nu}\delta F_{\nu\mu}=-\frac12F^{\mu\nu}\delta F_{\mu\nu}$, y
\[
\frac12F^{\mu\nu}\delta F_{\mu\nu}+\delta\left(\frac{1}{c\epsilon_0}
J^\mu A_\mu\right)+\partial_\nu\big(F^{\mu\nu}\delta A_\mu\big)=0.
\tag{111}
\]
Pero $\delta(F_{\mu\nu}F^{\mu\nu})=\delta F_{\mu\nu}F^{\mu\nu}+F_{\mu\nu}
\delta F^{\mu\nu}=2F^{\mu\nu}\delta F_{\mu\nu}$, es decir
\[
F^{\mu\nu}\delta F_{\mu\nu}=\frac12\delta\big(F_{\mu\nu}F^{\mu\nu}\big).
\tag{112}
\]
Introduciendo (112) en (111) e integrando en todo el espacio de
Minkowski $M^4$, con el teorema de Gauss para el término de borde,
\[
\delta\int_{M^4}\left(\frac14F_{\mu\nu}F^{\mu\nu}+\frac{1}{c\epsilon_0}
J^\mu A_\mu\right)d^4x-\int_{\partial M^4}F^{\nu\mu}\delta A_\mu\,dS_\nu
=0.
\]
Imponiendo la condición de borde $\delta A_\mu=0$ sobre $\partial M^4$,
\[
\delta\int_{M^4}\left(\frac14F_{\mu\nu}F^{\mu\nu}+\frac{1}{c\epsilon_0}
J^\mu A_\mu\right)d^4x=0. \tag{113}
\]
Es lógico postular la integral (114) como principio de acción del campo
electromagnético. Sin embargo, la integral debe respetar dos condiciones:
debe tener unidades de acción, $[S]=ML^2/T$ (115), y el término cinético
debe ser positivo. Las dimensiones de cada término son
$[F_{\mu\nu}F^{\mu\nu}]=[J^\mu A_\mu/(c\epsilon_0)]=M^2L^6/(T^4C^2)\neq
ML^2/T$, de modo que la integral no tiene unidades de acción. Para
repararlo multiplicamos por una constante dimensionada $\lambda$ (116):
$[\lambda]M^2L^6/(T^4C^2)=ML^2/T$, es decir
$[\lambda]=T^3C^2/(ML^4)$. Debemos construir $\lambda$ con las
constantes electromagnéticas $c$ y $\epsilon_0$ (pues
$\epsilon_0\mu_0=1/c^2$); por análisis dimensional,
\[
\lambda=\alpha\frac{\epsilon_0}{c}, \tag{117}
\]
con $\alpha$ adimensional, igual a $+1$ o $-1$ de modo que el término
cinético sea positivo. Entonces
\[
\int_{M^4}\left(\alpha\frac{\epsilon_0}{4c}F_{\mu\nu}F^{\mu\nu}+
\frac{\alpha}{c^2}J^\mu A_\mu\right)d^4x. \tag{118}
\]
Para hallar $\alpha$ analicemos el término cinético:
\[
F_{\mu\nu}F^{\mu\nu}=-2F_{0i}F_{0i}+F_{ij}F_{ij}=-2\left(\frac1c
\frac{\partial\vec A}{\partial t}-\vec\nabla\phi\right)^2+(\partial_iA_j-
\partial_jA_i)^2,
\]
que contiene el término cinético $-\dfrac{2}{c^2}\Big(\dfrac{\partial\vec
A}{\partial t}\Big)^2<0$. Por lo tanto se debe elegir $\alpha=-1$.

\begin{definicion}[Principio de acción de Maxwell]
\[
S_M=\int_{M^4}\left(-\frac{\epsilon_0}{4c}F_{\mu\nu}F^{\mu\nu}-\frac{1}
{c^2}J^\mu A_\mu\right)d^4x, \tag{119}
\]
con densidad Lagrangiana de Maxwell
\[
\mathscr{L}_M=-\frac{\epsilon_0}{4c}F_{\mu\nu}F^{\mu\nu}-\frac{1}{c^2}
J^\mu A_\mu . \tag{120}
\]
\end{definicion}

\section{Invariancia del principio de acción de Maxwell}

Como se ha mencionado, los principios de acción deben respetar los
postulados de la relatividad especial, es decir, deben ser invariantes
bajo transformaciones de Lorentz.

\begin{proposicionc}
El principio de acción de Maxwell (119) es invariante de Lorentz.
\end{proposicionc}

\begin{proof}
Sea la transformación de Lorentz $x'^\mu=\Lambda^\mu_{\ \nu}x^\nu$.
Entonces
\[
\begin{aligned}
S_M' &= \int_{M'^4}\left(-\frac{\epsilon_0}{4c}F'_{\mu\nu}F'^{\mu\nu}-
\frac{1}{c^2}J'^\mu A'_\mu\right)d^4x'\\
&= \int_{M^4}\left(-\frac{\epsilon_0}{4c}(\Lambda^{-1})^\alpha_{\ \mu}
(\Lambda^{-1})^\beta_{\ \nu}\Lambda^\mu_{\ \tau}\Lambda^\nu_{\ \rho}
F_{\alpha\beta}F^{\tau\rho}-\frac{1}{c^2}(\Lambda^{-1})^\alpha_{\ \mu}
\Lambda^\mu_{\ \beta}J^\beta A_\alpha\right)|\det\Lambda|\,d^4x\\
&= \int_{M^4}\left(-\frac{\epsilon_0}{4c}\delta^\alpha_\tau\delta^\beta_
\rho F_{\alpha\beta}F^{\tau\rho}-\frac{1}{c^2}\delta^\alpha_\beta J^\beta
A_\alpha\right)d^4x = S_M,
\end{aligned}
\]
usando $\det\Lambda=1$. \qedhere
\end{proof}

La invariancia de Lorentz significa que la cantidad $-\dfrac{\epsilon_0}
{4c}F_{\mu\nu}F^{\mu\nu}$ posee el mismo valor para todos los observadores
inerciales,
\[
-\frac{\epsilon_0}{4c}F_{\mu\nu}F^{\mu\nu}\Big|_K=-\frac{\epsilon_0}{4c}
F_{\mu\nu}F^{\mu\nu}\Big|_{K'}=-\frac{\epsilon_0}{4c}F_{\mu\nu}F^{\mu\nu}
\Big|_{K''}=\dots \tag{121}
\]
Para entender (121) expresémosla en términos de los campos:
\[
-\frac{\epsilon_0}{4c}F_{\mu\nu}F^{\mu\nu}=\frac{\epsilon_0}{2c}\Big(
\|\vec E\|^2-\|\vec H\|^2\Big). \tag{122}
\]
Evaluémosla en dos situaciones distintas.

\begin{ejemplo}[Condensador cargado]
Consideremos un condensador de placas infinitas con densidad de carga
$\sigma$, con los ejes de la figura.

\begin{center}
\begin{tikzpicture}[scale=1]
  \draw[very thick,colorprop] (0,1.6) -- (3.6,1.6);
  \draw[very thick,colordef] (0,0) -- (3.6,0);
  \node[left] at (0,1.6) {$+Q$};
  \node[left] at (0,0) {$-Q$};
  \foreach \x in {0.7,1.8,2.9} {
    \draw[-Stealth,thick,colorej] (\x,1.5) -- (\x,0.1);
  }
  \node[right] at (3.6,0.8) {$\vec E$};
  \draw[-Stealth] (4.6,0) -- (5.4,0) node[right] {$x$};
  \draw[-Stealth] (4.6,0) -- (4.6,0.8) node[above] {$z$};
\end{tikzpicture}
\end{center}

Con respecto a un S.R.I. $K$, el campo eléctrico y magnético son
\[
\vec E=\frac{\sigma}{\epsilon_0}\hat k, \qquad \vec H=\vec0. \tag{123}
\]
Evaluando (122),
\[
-\frac{\epsilon_0}{4c}F_{\mu\nu}F^{\mu\nu}=\frac{\epsilon_0}{2c}\left(
\Big\|\frac{\sigma}{\epsilon_0}\hat k\Big\|^2-0\right)=\frac{\sigma^2}
{2c\epsilon_0}>0. \tag{124}
\]
Sea ahora un S.R.I. $K'$ que se mueve con velocidad $\vec V=V\hat\jmath$
con respecto a $K$. De acuerdo con (80) o (81), los campos medidos por
$K'$ son
\[
\vec E\,'=\gamma\frac{\sigma}{\epsilon_0}\hat k, \qquad \vec H\,'=\gamma
\beta\frac{\sigma}{\epsilon_0}\hat\imath, \tag{125}
\]
y entonces
\[
-\frac{\epsilon_0}{4c}F'_{\mu\nu}F'^{\mu\nu}=\frac{\epsilon_0}{2c}\left(
\gamma^2\frac{\sigma^2}{\epsilon_0^2}-\gamma^2\beta^2\frac{\sigma^2}{
\epsilon_0^2}\right)=\frac{\sigma^2}{2c\epsilon_0}\gamma^2(1-\beta^2)=
\frac{\sigma^2}{2c\epsilon_0}>0. \tag{126}
\]
Se concluye que si un S.R.I. mide $-\frac{\epsilon_0}{4c}F_{\mu\nu}
F^{\mu\nu}>0$, entonces todos los S.R.I. medirán lo mismo.
\end{ejemplo}

\begin{ejemplo}[Solenoide]
Consideremos un solenoide muy largo, por el que circula una corriente
$I$ y que genera un campo magnético en su interior. Con respecto a un
S.R.I. $K$,
\[
\vec E=\vec0, \qquad \vec H=\frac{nI}{\epsilon_0c}\hat k, \tag{127}
\]
donde $n=N/L$, con $N$ el número de espiras. Evaluando (122),
\[
-\frac{\epsilon_0}{4c}F_{\mu\nu}F^{\mu\nu}=\frac{\epsilon_0}{2c}\left(0-
\Big\|\frac{nI}{\epsilon_0c}\hat k\Big\|^2\right)=-\frac{n^2I^2}{2c^3
\epsilon_0}<0. \tag{128}
\]
Para un S.R.I. $K'$ que se mueve con velocidad $\vec V=V\hat\jmath$, los
campos medidos son
\[
\vec E\,'=-\gamma\beta\frac{nI}{\epsilon_0c}\hat\imath, \qquad \vec H\,'=
\gamma\frac{nI}{\epsilon_0c}\hat k, \tag{129}
\]
y
\[
-\frac{\epsilon_0}{4c}F'_{\mu\nu}F'^{\mu\nu}=\frac{\epsilon_0}{2c}\left(
\gamma^2\beta^2\frac{n^2I^2}{\epsilon_0^2c^2}-\gamma^2\frac{n^2I^2}{
\epsilon_0^2c^2}\right)=-\frac{n^2I^2}{2c^3\epsilon_0}\gamma^2(1-\beta^2)
=-\frac{n^2I^2}{2c^3\epsilon_0}<0. \tag{130}
\]
Se concluye que si un S.R.I. mide $-\frac{\epsilon_0}{4c}F_{\mu\nu}
F^{\mu\nu}<0$, entonces todos los S.R.I. medirán lo mismo.
\end{ejemplo}

\section{Invariancia de gauge}

La teoría de Maxwell, como teoría física, debe respetar los postulados de
la relatividad especial: las ecuaciones de Maxwell deben ser tensores de
Lorentz y la integral variacional debe ser un escalar de Lorentz, y la
teoría de Maxwell respeta esta simetría. Sin embargo, posee una
\emph{segunda} simetría: una transformación de los cuadripotenciales
$A_\mu$ que deja las ecuaciones de campo invariantes en forma funcional,
la llamada \textbf{teoría de gauge}.

\begin{definicion}[Transformación de gauge]
Sea la transformación del cuadripotencial
\[
A'_\mu=A_\mu+\partial_\mu\alpha(x^\nu), \tag{131}
\]
donde $\alpha(x^\mu)$ es una función arbitraria de las coordenadas.
\end{definicion}

La transformación (131) actúa sobre los cuadripotenciales $A_\mu$ y no
sobre las coordenadas $x^\mu$: las coordenadas no cambian.

\begin{proposicionc}
La transformación de gauge (131) deja invariante el tensor de Faraday:
$F'_{\mu\nu}=F_{\mu\nu}$ (132).
\end{proposicionc}

\begin{proof}
\[
F'_{\mu\nu}=\partial_\mu(A_\nu+\partial_\nu\alpha)-\partial_\nu(A_\mu+
\partial_\mu\alpha)=\partial_\mu A_\nu-\partial_\nu A_\mu+\partial_\mu
\partial_\nu\alpha-\partial_\nu\partial_\mu\alpha=F_{\mu\nu}. \qedhere
\]
\end{proof}

Como consecuencia de (132), la primera ecuación de Maxwell queda
invariante,
\[
\partial_\nu F'^{\mu\nu}=-\frac{1}{c\epsilon_0}J^\mu\;\Longrightarrow\;
\partial_\nu F^{\mu\nu}=-\frac{1}{c\epsilon_0}J^\mu, \tag{133}
\]
donde $J^\mu$ no depende de $A_\mu$, solo de las coordenadas. La segunda
ecuación de Maxwell (68) se anula idénticamente al reemplazar $F_{\mu\nu}$
en función del cuadripotencial, así que no entra en la discusión. La
transformación (131) también deja invariante la ecuación (100):
\[
\Box A'_\mu-\partial_\mu(\partial_\nu A'^\nu)=\Box A_\mu-\partial_\mu
(\partial_\nu A^\nu)+\big(\Box\partial_\mu\alpha-\partial_\mu\Box\alpha
\big)=\frac{1}{c\epsilon_0}J_\mu .
\]
Esto implica que si $A_\mu$ es solución de las ecuaciones de Maxwell,
también $A_\mu+\partial_\mu\alpha$ lo es, y como $\alpha(x^\mu)$ es
arbitraria, las ecuaciones de Maxwell poseen infinitas soluciones: la
solución \textbf{no es única}.

\begin{ejemplo}[El solenoide]
Consideremos el solenoide con campo magnético $\vec H=\dfrac{nI}{\epsilon_0
c}\hat k$ (134); encontremos el potencial vectorial que satisface $\vec
H=-\vec\nabla\times\vec A$. Las ecuaciones a resolver son
\[
\frac{\partial A_y}{\partial x}-\frac{\partial A_x}{\partial y}=-
\frac{nI}{\epsilon_0c}, \qquad \frac{\partial A_z}{\partial y}-
\frac{\partial A_y}{\partial z}=0, \qquad \frac{\partial A_z}{\partial x}
-\frac{\partial A_x}{\partial z}=0. \tag{135}
\]
Como las soluciones son infinitas, mostramos solo algunas:
\[
\begin{aligned}
\vec A_1&=\frac{nIy}{\epsilon_0c}\hat\imath, \qquad \vec A_2=-\frac{nIx}
{\epsilon_0c}\hat\jmath, \qquad \vec A_3=\frac{nIy}{2\epsilon_0c}\hat\imath
-\frac{nIx}{2\epsilon_0c}\hat\jmath,\\
\vec A_4&=\frac{nIy}{\epsilon_0c}\hat\imath+\frac{nI}{\epsilon_0c}\cos
(kz)\hat k, \qquad \vec A_5=-\frac{nIx}{\epsilon_0c}\hat\jmath+\frac{nI}
{\epsilon_0c}e^{kz}\hat k,
\end{aligned} \tag{136}
\]
donde la cantidad $k$ es tal que $kz$ es adimensional. Estas soluciones
están relacionadas entre sí:
\[
\begin{aligned}
\vec A_2&=\vec A_1+\vec\nabla\!\left(-\frac{nIxy}{\epsilon_0c}\right),&
\vec A_3&=\vec A_2+\vec\nabla\!\left(\frac{nIxy}{2\epsilon_0c}\right),\\
\vec A_4&=\vec A_1+\vec\nabla\!\left(\frac{nI}{k\epsilon_0c}\sin(kz)
\right),& \vec A_5&=\vec A_2+\vec\nabla\!\left(\frac{nI}{k\epsilon_0c}
e^{kz}\right).
\end{aligned} \tag{137}
\]
\end{ejemplo}

Tener infinitas soluciones es un problema: si queremos analizar un sistema
electromagnético debemos tener un solo cuadripotencial. ¿Cuál debemos
usar? Cualquiera de ellas, pues todas son equivalentes: la relación
(131) es una \emph{relación de equivalencia}. Como todas son
equivalentes, escogemos la más simple, aquella que simplifique las
ecuaciones de campo, es decir, fijamos un solo $\alpha(x^\mu)$: esto se
llama \textbf{fijar el gauge}. Para simplificar la ecuación (100)
escogemos
\[
\partial_\mu A^\mu=0, \tag{138}
\]
llamada \textbf{gauge de Lorentz}. Introduciéndola en (100) se obtiene
\[
\Box A^\mu=\frac{1}{c\epsilon_0}J^\mu, \tag{139}
\]
que es una ecuación de onda con fuente.

\begin{ejemplo}
Apliquemos (138) a las soluciones (136) del solenoide (con $\phi=0$, de
modo que $\partial_\mu A^\mu=\vec\nabla\cdot\vec A$):
\[
\vec\nabla\cdot\vec A_1=\vec\nabla\cdot\vec A_2=\vec\nabla\cdot\vec A_3=0,
\qquad \vec\nabla\cdot\vec A_4=-\frac{nIk}{\epsilon_0c}\sin(kz)\neq0,
\qquad \vec\nabla\cdot\vec A_5=\frac{nIk}{\epsilon_0c}e^{kz}\neq0. \tag{140}
\]
\end{ejemplo}

Las soluciones $\vec A_4$ y $\vec A_5$ no satisfacen la fijación de gauge
(138). Esto se debe a la presencia de la función arbitraria $\alpha$:
si $A_\mu$ satisface (138), el cuadripotencial transformado no
necesariamente la satisface. En efecto,
\[
\partial^\mu A'_\mu=\partial^\mu(A_\mu+\partial_\mu\alpha)=\partial^\mu
A_\mu+\Box\alpha=\Box\alpha, \tag{141}
\]
y entonces $\Box\alpha\neq0\Rightarrow\partial^\mu A'_\mu\neq0$ (142).
La expresión (138) particiona el conjunto (infinito) de cuadripotenciales
que generan el campo electromagnético en dos subconjuntos disjuntos:
aquellos que satisfacen (138), y por tanto $\Box A^\mu=J^\mu/(c\epsilon_0)$,
y aquellos que no la satisfacen, y por tanto $\Box A^\mu-\partial^\mu
(\partial_\nu A^\nu)=J^\mu/(c\epsilon_0)$, lo que contradice la idea de
tener un solo cuadripotencial. Para salvar esta contradicción imponemos
la condición de fijación
\[
\Box\alpha=0, \tag{143}
\]
que implica que todos los cuadripotenciales satisfacen el gauge de
Lorentz (138). En el solenoide hay tres soluciones que lo satisfacen,
pero se elige la más simple, $\vec A_1$; las otras dos están relacionadas
con $\vec A_1$ por funciones $\alpha$ que satisfacen (143).

\begin{observacion}
La transformación de gauge (131) no está relacionada con una
transformación de Lorentz: las transformaciones de Lorentz relacionan
S.R.I.\ físicos, mientras que las transformaciones de gauge relacionan
cuadripotenciales independientes del observador.
\end{observacion}

\section{Teoría de Klein--Gordon}

Los físicos Oskar Klein y Walter Gordon propusieron en 1926 una ecuación
que describiera electrones, ecuación que lleva sus nombres. Sin embargo,
resultó que la ecuación de Klein--Gordon no describe electrones, sino
partículas de espín cero (cada componente de un campo de espín $1/2$
también la satisface, pero ese caso se describe con la ecuación de
Dirac). Describe correctamente a los \textbf{piones}, $\{\pi^+,\pi^0,
\pi^-\}$, partículas subatómicas portadoras de la interacción nuclear
entre los nucleones (protón y neutrón); en general describe partículas de
espín cero y compuestas por partículas más fundamentales.

La obtención de la ecuación de Klein--Gordon usa la relatividad especial
y la mecánica cuántica. Comencemos con la ecuación de la energía
relativista,
\[
E^2=c^2\|\vec p\|^2+m^2c^4, \tag{144}
\]
donde $E$ es la energía, $\vec p$ el momentum lineal tridimensional, $c$
la rapidez de la luz y $m$ la masa propia. Introducimos la prescripción
de operadores de la mecánica cuántica,
\[
\vec p\to\hat{\vec p}=-i\hbar\vec\nabla, \qquad E\to\hat E=i\hbar
\frac{\partial}{\partial t}, \tag{145}
\]
de modo que
\[
\hat E^2=c^2\|\hat{\vec p}\|^2+m^2c^4. \tag{146}
\]
Esta es una ecuación de operadores y por lo tanto debe actuar sobre una
función de onda $\phi(t,\vec x)$, llamada \textbf{campo escalar} o
\textbf{campo de Klein--Gordon}:
\[
-\hbar^2\frac{\partial^2\phi}{\partial t^2}=-c^2\hbar^2\nabla^2\phi+m^2c^4
\phi. \tag{147}
\]
Dividiendo por $c^2\hbar^2$ y usando notación relativista,
\[
-\partial_0^2\phi+\delta^{ij}\partial_i\partial_j\phi-\frac{m^2c^2}{\hbar^2}
\phi=0\;\Longrightarrow\;\eta^{\mu\nu}\partial_\mu\partial_\nu\phi-
\frac{m^2c^2}{\hbar^2}\phi=0,
\]
y así se obtiene la \textbf{ecuación de Klein--Gordon}
\[
\Box\phi-\frac{m^2c^2}{\hbar^2}\phi=0, \tag{148}
\]
donde $\Box$ es el operador d'Alembertiano.

Debemos ver si la ecuación de Klein--Gordon es un tensor de Lorentz bajo
$x'^\mu=\Lambda^\mu_{\ \nu}x^\nu$. Es directo observar que
\[
\Box'=\eta^{\mu\nu}\partial'_\mu\partial'_\nu=\eta^{\mu\nu}(\Lambda^{-1})^
\alpha_{\ \mu}(\Lambda^{-1})^\beta_{\ \nu}\partial_\alpha\partial_\beta=
\eta^{\alpha\beta}\partial_\alpha\partial_\beta=\Box, \tag{149}
\]
es decir, el d'Alembertiano es invariante de Lorentz. Solo falta saber
cómo transforma el campo escalar. El campo escalar no tiene índices
libres, por lo que debe ser un escalar de Lorentz para que la ecuación
de Klein--Gordon sea un tensor de Lorentz.

\begin{definicion}[Campo escalar de Klein--Gordon]
El campo escalar de Klein--Gordon es un escalar de Lorentz:
\[
\phi'(x'^\alpha)=\phi(x^\beta), \tag{150}
\]
es decir, es un invariante de Lorentz.
\end{definicion}

Se concluye que la ecuación de Klein--Gordon es un escalar de Lorentz:
\[
\Box'\phi'(x'^\alpha)-\frac{m^2c^2}{\hbar^2}\phi'(x'^\alpha)=\Box\phi(x^
\beta)-\frac{m^2c^2}{\hbar^2}\phi(x^\beta)=0. \tag{151}
\]

Usando ingeniería inversa para obtener el principio de acción,
multiplicamos la ecuación de Klein--Gordon por $\delta\phi$:
\[
\Box\phi\,\delta\phi-\frac{m^2c^2}{\hbar^2}\phi\,\delta\phi=\partial_\mu
\big(\eta^{\mu\nu}\partial_\nu\phi\big)\delta\phi-\delta\left(\frac{m^2c^2}
{2\hbar^2}\phi^2\right)
\]
\[
=-\eta^{\mu\nu}\partial_\nu\phi\,\delta(\partial_\mu\phi)-\delta\left(
\frac{m^2c^2}{2\hbar^2}\phi^2\right)+\partial_\mu\big(\eta^{\mu\nu}
\partial_\nu\phi\,\delta\phi\big). \tag{152}
\]
Pero $\delta(\eta^{\mu\nu}\partial_\mu\phi\,\partial_\nu\phi)=2\eta^{\mu\nu}
\partial_\nu\phi\,\delta(\partial_\mu\phi)$, es decir
\[
\eta^{\mu\nu}\partial_\nu\phi\,\delta(\partial_\mu\phi)=\frac12\delta
\big(\eta^{\mu\nu}\partial_\mu\phi\,\partial_\nu\phi\big). \tag{153}
\]
Introduciendo (153) en (152),
\[
\delta\left(-\frac12\eta^{\mu\nu}\partial_\mu\phi\,\partial_\nu\phi-
\frac{m^2c^2}{2\hbar^2}\phi^2\right)+\partial_\mu\big(\eta^{\mu\nu}
\partial_\nu\phi\,\delta\phi\big)=0. \tag{154}
\]
Integrando en $M^4$ y usando el teorema de Gauss en el término de
borde, con la condición de borde $\delta\phi=0$ sobre $\partial M^4$,
\[
\delta\int_{M^4}\left(-\frac12\eta^{\mu\nu}\partial_\mu\phi\,\partial_\nu
\phi-\frac{m^2c^2}{2\hbar^2}\phi^2\right)d^4x=0. \tag{155}
\]
Lo lógico sería llamar a esta integral el principio de acción del campo
escalar, pero debe tener unidades de acción, $ML^2T^{-1}$. A diferencia
de Maxwell, multiplicar por una constante dimensionada $\lambda$ no
sirve: definiendo $\phi'=\sqrt\lambda\,\phi$ la constante desaparecería,
pues el campo aparece al cuadrado. Por lo tanto, no es necesario
multiplicar por una constante dimensionada; en cambio, las unidades del
campo se deducen exigiendo que $S$ tenga unidades de acción:
\[
\left[\frac{m^2c^2}{\hbar^2}\phi^2d^4x\right]=\frac{ML^2}{T}
\;\Longrightarrow\; \frac{1}{L^2}[\phi]^2L^4=\frac{ML^2}{T}
\;\Longrightarrow\; [\phi]=\frac{M^{1/2}}{T^{1/2}}. \tag{156}
\]
Verifiquemos que el término cinético sea positivo:
\[
-\frac12\eta^{\mu\nu}\partial_\mu\phi\,\partial_\nu\phi=\frac{1}{2c^2}
\dot\phi^2-\frac12\partial_i\phi\,\partial_i\phi,
\]
es positivo, por lo que no hay que cambiar nada.

\begin{definicion}[Principio de acción del campo escalar libre]
\[
S=\int_{M^4}\left(-\frac12\eta^{\mu\nu}\frac{\partial\phi}{\partial x^\mu}
\frac{\partial\phi}{\partial x^\nu}-\frac{m^2c^2}{2\hbar^2}\phi^2\right)
d^4x. \tag{157}
\]
\end{definicion}

\section{Campo escalar de Klein--Gordon con fuente}

La ecuación de Klein--Gordon con fuente describe un campo escalar libre
en presencia de fuentes externas, y se escribe
\[
\Box\phi-\frac{m^2c^2}{\hbar^2}\phi=-J, \tag{158}
\]
donde $J(x^\mu)$ es la densidad de fuente. Se diferencia de la ecuación
sin fuente ($J=0$) porque incluye la interacción con una fuente externa,
como campos electromagnéticos o partículas. El principio de acción para
este caso es
\[
S=\int_{M^4}\left(-\frac12\eta^{\mu\nu}\frac{\partial\phi}{\partial x^\mu}
\frac{\partial\phi}{\partial x^\nu}-\frac{m^2c^2}{2\hbar^2}\phi^2+J\phi
\right)d^4x. \tag{159}
\]
